\chapter{Problema}
\label{cap:problema}

A otimização da produção de petróleo em poços é uma tarefa crítica de engenharia, cuja eficiência está diretamente ligada à acurácia da previsão da Vazão Diária de Petróleo ($\hat{y}$) em função das variáveis operacionais e de reservatório ($X$). As complexas interações entre estas variáveis frequentemente dificultam a obtenção de estimativas precisas por métodos tradicionais, que falham em mapear a função $f(X)$ de forma satisfatória. Tais iterações entre as variáveis, pode classificar o problema quanto à sua linearidade, se este será \textbf{linear} ou \textbf{não-linear}, e dada a grande quantidade de variáveis, é muito difícil classificar o problema apenas com a observação dos dados. Sendo assim, o problema abordado será modelado de ambas as formas e serão utilizadas as métricas dos modelos de \gls{ML}, apresentados no Capítulo \ref{cap:experimentos}, para determinar a natureza do problema, levando em consideração que foram utilizados modelos diferentes, onde alguns assumem o problema como linear e outros como não-linear, o que será melhor detalhado no Capítulo \ref{cap:proposta-e-metodologia}.

Abaixo será apresentada a formulação matemática do problema, bem como os 2 tipos de modelagem do problema de acordo com as naturezas dos algoritmos, e por fim do problema de otimização das métricas.

\section{Formulação do Problema}
\label{subsec:formulacao}

Com o entendimento do problema e de suas variáveis, foi possível desenvolver uma possível formulação matemática, de modo a tornar possível a utilização de algoritmos de otimização, e encontrar a solução ótima. Abaixo está a modelagem desenvolvida:

\begin{align}
    \max_{\mathbf{X}} & \quad \hat{y} = f(\mathbf{X}) \label{general_fo} \\
    \text{s.a.} & \quad L_{\min, j} \le x_j \le L_{\max, j}, \quad \forall j \in \{1, 2, \dots, p\} \label{general_eq1} \\
    & \quad x_{\text{Grau\_API}} = C_{\text{API}} \label{general_eq2} \\
    & \quad x_{\text{Idade}} = C_{\text{Idade}} \label{general_eq3} \\
    & \quad x_{\text{PCS}} = C_{\text{PCS}} \label{general_eq4} \\
    & \quad x_{\text{Densidade GLP Gás}} = C_{\text{Densidade GLP Gás}} \label{general_eq5} \\
    & \quad x_{\text{Densidade GLP Líquido}} = C_{\text{Densidade GLP Líquido}} \label{general_eq6} \\
    & \quad x_{\text{Lâmina D'Água}} = C_{\text{Lâmina D'Água}} \label{general_eq7} \\
    & \quad x_{\text{Distância Poço k}} = C_{\text{Distância Poço k}} \label{general_eq8} \\
    & \quad x_{\text{Poços no Raio de 2 km}} = C_{\text{Poços no Raio de 2 km}} \label{general_eq9} \\
    & \quad x_j \in \mathbb{R}_{+}, \quad \forall j \in \{1, 2, \dots, p\} \label{general_eq10}
\end{align}

Nesta modelagem, foi definida na equação (\ref{general_fo}) a função objetivo, que estabelece a meta de maximizar a produção estimada $\hat{y}$. A função $f(\mathbf{X})$ consiste na relação preditiva aprendida pelos algoritmos de \gls{ML}, como \textit{Linear Regression Forest} ou \textit{LightGBM}, capaz de capturar interações entre as variáveis, sejam essas lineares ou não lineares. A definição da função objetivo apresentada veio com base no escopo desse trabalho, que busca utilizar tais modelos para encontrar aquele que melhor se adapte aos dados apresentados.

As variáveis deste modelo são representadas pelo vetor $\mathbf{X}=\{x_1, x_2, \dots, x_p\}$, composto pelos parâmetros operacionais que podem ser manipulados para influenciar a vazão. O domínio definido para essas variáveis na equação (\ref{general_eq10}) como reais não negativos ($\mathbb{R}_{+}$) reflete a natureza física de grandezas como vazões e intervalos temporais.

A equação \ref{general_eq1} determina o intervalo das variáveis de decisão $x_j \in \mathbf{X}$. 
Tais variáveis devem respeitar os limites mínimos ($L_{\min}$) e máximos ($L_{\max}$) de segurança e capacidade dos equipamentos.

Por fim, as equações (\ref{general_eq2})-(\ref{general_eq9}) determinam os parâmetros do problema, de modo que tais informações são relacionadas à um poço específico, e não podem ser alteradas.

\section{Problema Linear}

Os modelos lineares supervisionados assumem uma relação funcional paramétrica e aditiva entre uma matriz de atributos preditores e uma variável resposta contínua. O objetivo consiste em estimar o vetor de coeficientes $\boldsymbol{\beta} \in \mathbb{R}^{p+1}$ que minimiza a soma dos \gls{RSS}, acrescida ou não de um termo de regularização para conter a variância das estimativas. A formulação geral de otimização irrestrita para essa classe de regressores é expressa por:

\begin{align}
\min_{\boldsymbol{\beta} \in \mathbb{R}^{p+1}} \quad J(\boldsymbol{\beta}) = \sum_{i=1}^{n} \left(y_i - \left(\beta_0 + \sum_{j=1}^{p} \beta_j x_{ij}\right)\right)^2 + \lambda \cdot R(\boldsymbol{\beta}) \label{eq:fo_linear_geral}
\end{align}

\noindent em que $y_i \in \mathbb{R}$ representa o valor escalar observado da variável resposta para a $i$-ésima instância ($i = 1, \dots, n$), e $x_{ij} \in \mathbb{R}$ denota o valor associado ao $j$-ésimo atributo preditivo ($j = 1, \dots, p$). O vetor de parâmetros a ser estimado decompõe-se no termo independente ou intercepto ($\beta_0 \in \mathbb{R}$) e nos pesos lineares associados a cada dimensão ($\beta_j \in \mathbb{R}$). O parâmetro $\lambda \ge 0$ atua como o hiperparâmetro de regularização, ditando a intensidade da penalidade sobre a magnitude dos coeficientes.

A partir da estrutura apresentada na equação \ref{eq:fo_linear_geral}, derivam-se três variantes fundamentais de regressão linear em função da definição adotada para o funcional $R(\boldsymbol{\beta})$:

\begin{itemize}
    \item \textbf{Regressão Linear por \gls{OLS}:} Obtida definindo $\lambda = 0$. Sob essa condição, o ajuste depende estritamente da minimização dos desvios quadráticos amostrais. Trata-se de um estimador não tendencioso de variância mínima sob as hipóteses do Teorema de Gauss-Markov, porém suscetível à instabilidade numérica e sobreajuste na presença de multicolinearidade.
    
    \item \textbf{Regressão Ridge (Regularização $L\_2$):} Configura-se aplicando a norma euclidiana quadrática dos pesos preditores, isto é, $R(\boldsymbol{\beta}) = \Vert{}\boldsymbol{\beta}_{1:p}\Vert{}_2^2 = \sum_{j=1}^{p} \beta_j^2$. A penalização quadrática contrai uniformemente os coeficientes em direção a zero, mitigando a variância do estimador frente a variáveis altamente correlacionadas, embora não realize a seleção esparsa de variáveis ($\beta_j \neq 0, \forall j$).

    \item \textbf{Regressão Lasso (Regularização $L\_1$):} Estrutura-se a partir da penalização pela norma de Manhattan, definida por $R(\boldsymbol{\beta}) = \Vert{}\boldsymbol{\beta}_{1:p}\Vert{}_1 = \sum_{j=1}^{p} \vert{}\beta_j\vert{}$. Devido à geometria dos contornos da norma $L_1$, que contém vértices nos eixos das coordenadas, essa formulação força coeficientes de variáveis irrelevantes ou redundantes a assumirem exatamente o valor zero para valores adequados de $\lambda$, atuando concomitantemente como estimador preditivo e mecanismo de seleção de atributos (\textit{feature selection}).
\end{itemize}

Ressalta-se que o termo independente $\beta_0$ é comumente excluído da função de penalidade $R(\boldsymbol{\beta})$, assegurando que a penalização incida exclusivamente sobre as inclinações relativas às variáveis explicativas e preserve o ajuste médio da distribuição amostral.

\section{Problema Não Linear}

A fim de superar a rigidez funcional das abordagens lineares e modelar interações mais complexas entre os atributos, é necessário recorrer a modelos de aprendizado supervisionado não linear. De forma geral, esses algoritmos assumem que a variável contínua observada $y_i \in \mathbb{R}$ relaciona-se com o vetor de atributos preditores $\mathbf{x}_i \in \mathbb{R}^p$ por meio de uma função arbitrária $f: \mathbb{R}^p \to \mathbb{R}$ sujeita a ruído estocástico:

\begin{align}
y_i = f(\mathbf{x}_i) + \varepsilon_i, \quad \text{com } \mathbb{E}[\varepsilon_i \mid \mathbf{x}_i] = 0 \text{ e } \operatorname{Var}(\varepsilon_i \mid \mathbf{x}_i) = \sigma^2 \label{eq:regressao_nao_linear_geral}
\end{align}

Em vez de restringir a busca a um hiperplano paramétrico estrito, o objetivo consiste em aproximar a função desconhecida $f(\cdot)$ dentro de um espaço de hipóteses $\mathcal{F}$, minimizando o risco empírico regularizado:

\begin{align}
\min_{f \in \mathcal{F}} \quad \frac{1}{n} \sum_{i=1}^{n} \mathcal{L}(y_i, f(\mathbf{x}_i)) + \Omega(f) \label{eq:otimizacao_nao_linear}
\end{align}

\noindent em que $L(\cdot)$ representa a função de perda (como o erro quadrático ou a perda insensível) e $\Omega(f)$ impõe uma penalidade sobre a complexidade funcional para mitigar o sobreajuste.

Nesta pesquisa, foram avaliados cinco modelos não lineares, distribuídos em três famílias metodológicas de \gls{ML}: algoritmos baseados em instâncias e vizinhança, algoritmos baseados em \gls{SVM}, e algoritmos baseados em árvores de decisão.

\subsection{Modelos baseados em Instâncias e Vizinhança}

A família de modelos baseados em instâncias e vizinhança fundamenta-se no princípio do aprendizado preguiçoso (\textit{lazy learning}), no qual não há uma fase de calibração de parâmetros globais ou estimação prévia de hiperplanos. O algoritmo investigado nesta categoria foi o \gls{KNN}, amplamente consolidado em problemas de classificação e adaptado para a regressão contínua \cite{ibm_knn}.

Dada uma nova observação $\mathbf{x}^* \in \mathbb{R}^p$, o modelo estima a variável de resposta contínua calculando a média aritmética dos valores observados entre os $k$ vizinhos mais próximos no espaço vetorial de atributos:

\begin{equation}
\hat{f}_{\text{KNN}}(\mathbf{x}^*) = \frac{1}{k} \sum_{i \in \mathcal{N}_k(\mathbf{x}^*)} y_i \label{eq:predicao_knn}
\end{equation}

\noindent em que $\mathcal{N}_k(\mathbf{x}^*)$ denota o conjunto de índices das $k$ instâncias amostrais com menor distância em relação a $\mathbf{x}^*$. A proximidade métrica entre as instâncias no espaço $\mathbb{R}^p$ pode ser generalizada pela distância de Minkowski com parâmetro $q \ge 1$:

\begin{equation}
d_q(\mathbf{x}^*, \mathbf{x}_i) = \left( \sum_{j=1}^{p} |x^*_j - x_{ij}|^q \right)^{\frac{1}{q}} \label{eq:distancia_minkowski}
\end{equation}

A partir da formulação geral em (\ref{eq:distancia_minkowski}), derivam-se métricas específicas em função do valor atribuído a $q$:
\begin{itemize}
    \item \textbf{Distância Manhattan ($q = 1$):} baseada na norma $\ell_1$, soma as diferenças absolutas em cada dimensão:
    \begin{equation}
    d_{\text{Manhattan}}(\mathbf{x}^*, \mathbf{x}_i) = \sum_{j=1}^{p} |x^*_j - x_{ij}| \label{eq:distancia_manhattan}
    \end{equation}
    
    \item \textbf{Distância Euclidiana ($q = 2$):} correspondente à norma euclidiana padrão $\ell_2$, adotada como configuração padrão neste trabalho:
    \begin{equation}
    d_{\text{Euclidiana}}(\mathbf{x}^*, \mathbf{x}_i) = \sqrt{\sum_{j=1}^{p} (x^*_j - x_{ij})^2} \label{eq:distancia_euclidiana}
    \end{equation}
\end{itemize}

A não linearidade do estimador manifesta-se pela partição local do domínio contínuo em células de Voronoi implícitas, permitindo que a resposta se adapte a superfícies de produção complexas sem hipóteses restritivas sobre a distribuição dos dados. Ressalta-se que, devido à sensibilidade intrínseca das métricas de distância em (\ref{eq:distancia_minkowski}) à magnitude dos atributos preditores, a padronização das variáveis via \texttt{StandardScaler} no \texttt{Pipeline} foi estritamente indispensável para assegurar pesos homogêneos a todas as dimensões durante a busca dos vizinhos mais próximos.

\subsection{Algoritmos baseados em \gls{SVM}}

O grupo de algoritmos baseados em \gls{SVM} busca assegurar alta capacidade de generalização a partir de um número limitado de padrões de treinamento, fundamentando-se no princípio da \gls{SRM}. Tal princípio atua equilibrando simultaneamente a minimização do risco empírico e o controle da dimensão \gls{VC}, a qual expressa a capacidade de representação do modelo por meio do tamanho máximo de um conjunto de dados passível de ser perfeitamente separado (ou rotulado) pela família de funções admitida \cite{basak_svr}. 

O conceito de \gls{SVM} foi formalizado por \citeonline{vapnik1995nature}, estabelecendo que o estimador projeta o vetor de atributos de entrada $\mathbf{x}$ em um espaço de características de alta dimensionalidade $\mathcal{F}$ por meio de um mapeamento não linear a priori, ambiente no qual se constrói um hiperplano ótimo de separação que maximiza a margem geométrica entre os padrões.

No trabalho em questão, utilizou-se o algoritmo \gls{SVR}, que estende a formulação original de classificação para a predição de variáveis contínuas, adotando-se especificamente a sua formulação não linear. Dada uma base de treinamento com $l$ instâncias $\{(\mathbf{x}_1, y_1), \dots, (\mathbf{x}_l, y_l)\}$, onde $\mathbf{x}_i \in \mathbb{R}^p$ e $y_i \in \mathbb{R}$, o objetivo consiste em determinar uma função linear no espaço transformado, $f(\mathbf{x}) = \langle \boldsymbol{\omega}, \Phi(\mathbf{x}) \rangle + b$, que apresente desvio máximo de $\varepsilon$ em relação aos alvos observados $y_i$ e mantenha a máxima regularidade possível (minimizando $\|\boldsymbol{\omega}\|^2$), em que $\langle \cdot, \cdot \rangle_{\mathcal{F}}$ denota o produto escalar no espaço de características $\mathcal{F}$. Para tanto, adota-se a função de perda $\varepsilon$-insensível de Vapnik:

\begin{equation}
    |\xi|_\varepsilon = 
    \begin{cases}
        0, & \text{se } |\xi| < \varepsilon \\
        |\xi| - \varepsilon, & \text{caso contrário}
    \end{cases} 
    \label{eq:perda_epsilon}
\end{equation}

\noindent em que desvios contidos no intervalo $[-\varepsilon, +\varepsilon]$ não recebem penalização, conferindo imunidade a variações estocásticas nas séries temporais de vazão \cite{basak_svr}.

Para acomodar desvios que eventualmente ultrapassem os limites de tolerância do tubo, introduzem-se variáveis de folga $\xi_i$ e $\xi_i^*$ para erros positivos e negativos, formulando o problema de otimização primal convexo:

\begin{align}
    \min_{\boldsymbol{\omega}, b, \boldsymbol{\xi}, \boldsymbol{\xi}^*} & \quad \frac{1}{2}\|\boldsymbol{\omega}\|^2 + C \sum_{i=1}^{l} (\xi_i + \xi_i^*) \label{eq:primal_svr}\\
    \text{s.a.} & \quad y_i - \langle \boldsymbol{\omega}, \Phi(\mathbf{x}_i) \rangle - b \le \varepsilon + \xi_i, \quad \forall i \in \{1, \dots, l\} \label{eq:restr_primal1} \\
    & \quad \langle \boldsymbol{\omega}, \Phi(\mathbf{x}_i) \rangle + b - y_i \le \varepsilon + \xi_i^*, \quad \forall i \in \{1, \dots, l\} \label{eq:restr_primal2} \\
    & \quad \xi_i, \xi_i^* \ge 0, \quad \forall i \in \{1, \dots, l\} \label{eq:restr_primal3}
\end{align}

\noindent onde a constante de regularização $C > 0$ estabelece o \textit{trade-off} entre a regularidade de $f$ e a tolerância aos desvios marginais.

A função de mapeamento não linear $\Phi: \mathbb{R}^p \to \mathcal{F}$ desempenha o papel central de projetar os vetores de atributos originais do espaço de entrada para um espaço de características de dimensão potencialmente infinita $\mathcal{F}$. A fundamentação teórica dessa transformação repousa no fato de que relações complexas e intrinsecamente não lineares no espaço original tornam-se passíveis de ajuste linear quando representadas nessa dimensionalidade superior \cite{muller2001introduction}. Contudo, a avaliação explícita das coordenadas de $\Phi(\mathbf{x})$ é frequentemente inviável do ponto de vista analítico e computacional. Essa barreira é contornada pelo artifício do núcleo, no qual o produto interno $\langle \Phi(\mathbf{x}_i), \Phi(\mathbf{x}_j) \rangle_{\mathcal{F}}$ é calculado de maneira estritamente implícita no espaço original por meio da função de núcleo $k(\mathbf{x}_i, \mathbf{x}_j)$, dispensando o conhecimento direto da forma funcional exata de $\Phi(\cdot)$.

A resolução eficiente em alta dimensão fundamenta-se na formulação dual de Lagrange. Introduzindo os multiplicadores $\alpha_i, \alpha_i^* \ge 0$ (associados ao tubo de tolerância) e $\eta_i, \eta_i^* \ge 0$ (associados às folgas $\xi_i, \xi_i^*$), a função Lagrangiana $L$ expressa-se por:

\begin{align}
    L = & \quad \frac{1}{2}\|\boldsymbol{\omega}\|^2 + C\sum_{i=1}^{l}(\xi_i + \xi_i^*) - \sum_{i=1}^{l}\alpha_i(\varepsilon + \xi_i - y_i + \langle \boldsymbol{\omega}, \Phi(\mathbf{x}_i) \rangle + b) \nonumber \\
    & - \sum_{i=1}^{l}\alpha_i^*(\varepsilon + \xi_i^* + y_i - \langle \boldsymbol{\omega}, \Phi(\mathbf{x}_i) \rangle - b) - \sum_{i=1}^{l}(\eta_i\xi_i + \eta_i^*\xi_i^*) \label{eq:lagrangiano_svr}
\end{align}

Pelas condições de ponto de sela para otimalidade, a anulação das derivadas parciais de $L$ em relação às variáveis primais $(\boldsymbol{\omega}, b, \xi_i, \xi_i^*)$ resulta em:

\begin{align}
    \frac{\partial L}{\partial b} = \sum_{i=1}^{l}(\alpha_i^* - \alpha_i) = 0 \label{eq:kkt_deriv_b} \\
    \frac{\partial L}{\partial \boldsymbol{\omega}} = \boldsymbol{\omega} - \sum_{i=1}^{l}(\alpha_i - \alpha_i^*)\Phi(\mathbf{x}_i) = 0 \quad \Longrightarrow \quad \boldsymbol{\omega} = \sum_{i=1}^{l}(\alpha_i - \alpha_i^*)\Phi(\mathbf{x}_i) \label{eq:vetor_omega} \\
    \frac{\partial L}{\partial \xi_i^{(*)}} = C - \alpha_i^{(*)} - \eta_i^{(*)} = 0 \label{eq:kkt_deriv_xi}
\end{align}

Substituindo (\ref{eq:kkt_deriv_b}), (\ref{eq:vetor_omega}) e (\ref{eq:kkt_deriv_xi}) na Equação (\ref{eq:lagrangiano_svr}), obtém-se o problema quadrático dual:

\begin{align}
    \max_{\boldsymbol{\alpha}, \boldsymbol{\alpha}^*} & \quad -\frac{1}{2}\sum_{i=1}^{l}\sum_{j=1}^{l}(\alpha_i - \alpha_i^*)(\alpha_j - \alpha_j^*) k(\mathbf{x}_i, \mathbf{x}_j) - \varepsilon\sum_{i=1}^{l}(\alpha_i + \alpha_i^*) + \sum_{i=1}^{l}y_i(\alpha_i - \alpha_i^*) \label{eq:dual_svr} \\
    \text{s.a.} & \quad \sum_{i=1}^{l}(\alpha_i - \alpha_i^*) = 0 \label{eq:restr_dual1} \\
    & \quad 0 \le \alpha_i, \alpha_i^* \le C, \quad \forall i \in \{1, \dots, l\} \label{eq:restr_dual2}
\end{align}

A função de predição final resulta da expansão por vetores de suporte expressa em (\ref{eq:vetor_omega}):

\begin{equation}
    \hat{f}_{\text{SVR}}(\mathbf{x}^*) = \sum_{i=1}^{l} (\alpha_i - \alpha_i^*) k(\mathbf{x}_i, \mathbf{x}^*) + b \label{eq:predicao_final_svr}
\end{equation}

Pelas condições de folga complementar de \gls{KKT}, os produtos entre variáveis duais e restrições primais anulam-se no ponto ótimo:

\begin{align}
    \alpha_i(\varepsilon + \xi_i - y_i + \langle \boldsymbol{\omega}, \Phi(\mathbf{x}_i) \rangle + b) &= 0 \label{eq:kkt_cond1} \\
    \alpha_i^*(\varepsilon + \xi_i^* + y_i - \langle \boldsymbol{\omega}, \Phi(\mathbf{x}_i) \rangle - b) &= 0 \label{eq:kkt_cond2} \\
    (C - \alpha_i)\xi_i = 0 \quad \text{e} \quad (C - \alpha_i^*)\xi_i^* &= 0 \label{eq:kkt_cond3}
\end{align}

Dentre as propriedades analíticas destacadas por \citeonline{basak_svr}, sobressaem-se:
\begin{enumerate}
    \item \textbf{Esparsidade (\textit{Sparsity}):} para amostras contidas no interior do tubo ($|f(\mathbf{x}_i) - y_i| < \varepsilon$), tem-se $\xi_i = \xi_i^* = 0$, implicando necessariamente $\alpha_i = \alpha_i^* = 0$. Apenas observações com resíduo igual ou superior a $\varepsilon$ apresentam coeficientes não nulos, sendo denominadas \textit{vetores de suporte};
    \item \textbf{Exclusão mútua:} como uma observação não pode violar simultaneamente as duas margens, decorre que $\alpha_i \cdot \alpha_i^* = 0$;
    \item \textbf{Cálculo do termo de viés ($b$):} para vetores de suporte situados estritamente na margem livre ($0 < \alpha_i^{(*)} < C$), tem-se $\xi_i^{(*)} = 0$, permitindo isolar $b$ a partir de (\ref{eq:kkt_cond1}) e (\ref{eq:kkt_cond2}):
    \begin{equation}
    b = 
    \begin{cases}
    y_i - \sum_{j=1}^{l}(\alpha_j - \alpha_j^*) k(\mathbf{x}_j, \mathbf{x}_i) - \varepsilon, & \text{para } \alpha_i \in (0, C) \\
    y_i - \sum_{j=1}^{l}(\alpha_j - \alpha_j^*) k(\mathbf{x}_j, \mathbf{x}_i) + \varepsilon, & \text{para } \alpha_i^* \in (0, C)
    \end{cases} \label{eq:calculo_bias}
    \end{equation}
\end{enumerate}

A escolha da função de núcleo (\textit{kernel}) $k(\mathbf{x}_i, \mathbf{x}_j)$ determina a geometria do mapeamento implícito e a classe funcional do estimador. No presente trabalho, o espaço de busca hiperparamétrica avaliou quatro formulações mercerianas clássicas:

\begin{enumerate}
    \item \textbf{Linear:}
    \begin{equation}
        k(\mathbf{x}_i, \mathbf{x}_j) = \langle \mathbf{x}_i, \mathbf{x}_j \rangle \label{eq:kernel_linear}
    \end{equation}
    
    \item \textbf{\gls{RBF}:}
    \begin{equation}
        k(\mathbf{x}_i, \mathbf{x}_j) = \exp\left(-\gamma \|\mathbf{x}_i - \mathbf{x}_j\|_2^2\right), \quad \gamma > 0 \label{eq:kernel_rbf}
    \end{equation}
    
    \item \textbf{Polinomial:}
    \begin{equation}
        k(\mathbf{x}_i, \mathbf{x}_j) = \left(\gamma \langle \mathbf{x}_i, \mathbf{x}_j \rangle + r\right)^d, \quad \gamma > 0, \; d \in \mathbb{N}^* \label{eq:kernel_poly}
    \end{equation}
    
    \item \textbf{Sigmoide:}
    \begin{equation}
        k(\mathbf{x}_i, \mathbf{x}_j) = \tanh\left(\gamma \langle \mathbf{x}_i, \mathbf{x}_j \rangle + r\right) \label{eq:kernel_sigmoid}
    \end{equation}
\end{enumerate}

\noindent em que $\gamma$ governa a largura da influência local ou escala das projeções, $r$ denota o termo independente de deslocamento e $d$ representa o grau polinomial. Todos os hiperparâmetros estruturais ($C$ e $\varepsilon$), bem como os parâmetros intrínsecos de cada família de núcleos ($\gamma$, $r$ e $d$), foram sistematicamente calibrados por validação cruzada por meio de \texttt{GridSearchCV}. 

Ademais, em virtude da sensibilidade às magnitudes relativas dos preditores — evidenciada tanto pela norma euclidiana presente em (\ref{eq:kernel_rbf}) quanto pelos produtos escalares sucessivos em (\ref{eq:kernel_linear}), (\ref{eq:kernel_poly}) e (\ref{eq:kernel_sigmoid}) —, a padronização das covariáveis via \texttt{StandardScaler} no fluxo do \texttt{Pipeline} foi mantida estritamente para todos os núcleos avaliados, assegurando estabilidade numérica e homogeneidade dimensional durante a otimização quadrática.

\subsection{Algoritmos baseados em Árvores de Decisão}

As árvores de decisão para regressão fundamentam-se no particionamento binário recursivo e hierárquico do espaço de atributos em sub-regiões ortogonais disjuntas, associando a cada nó terminal (folha) um valor constante determinado pela média amostral local das respostas. Conforme revisado historicamente por \citeonline{loh2014fifty}, a metodologia consolidou-se a partir de pioneiros como o \gls{AID} e, posteriormente, o \gls{CART}, nos quais a escolha da melhor divisão em cada nó decorre da busca exaustiva pelo atributo e ponto de corte que maximizem a redução da soma dos quadrados dos resíduos locais. Não obstante a sua elevada interpretabilidade estrutural e capacidade intrínseca de lidar com não linearidades, estimadores baseados em árvores individuais apresentam alta variância amostral e instabilidade numérica frente a pequenas perturbações nos dados de treinamento, fatores que impulsionaram historicamente a sua consolidação como estimadores-base em arquiteturas de comitês (\textit{ensembles}).

No escopo deste trabalho, foram avaliados os seguintes algoritmos baseados em árvores: Random Forest, Gradient Boosting.

\subsubsection{Random Forest}

O algoritmo \textit{Random Forest} fundamenta-se na extensão do paradigma de \textit{bagging} aplicado a árvores de decisão, tendo sido formalizado originalmente por \citeonline{Breiman_2001_RandomForests}. O modelo opera construindo uma coleção de estimadores $\{h(\mathbf{x}, \Theta_k), k = 1, \dots, K\}$, denominados árvores de decisão, em que cada unidade é ajustada sobre uma partição de dados gerada por reamostragem com reposição (\textit{bootstrap}) e os vetores aleatórios $\Theta_k$ determinam a amostragem independente de um subconjunto restrito de atributos em cada divisão de nó.

Considerando um problema de regressão em que as respostas observadas $Y \in \mathbb{R}$ são variáveis contínuas associadas a atributos $\mathbf{X} \in \mathbb{R}^p$, com amostra de treinamento $\mathcal{D}_n = \{(\mathbf{X}_1, Y_1), \dots, (\mathbf{X}_n, Y_n)\}$ sorteada independentemente a partir da distribuição conjunta $(\mathbf{X}, Y)$, o erro quadrático médio de generalização para qualquer estimador $h(\mathbf{X})$ é definido formalmente por:

\begin{equation}
    \text{PE}^*(h) = \mathbb{E}_{\mathbf{X}, Y}\left[\left(Y - h(\mathbf{X})\right)^2\right] \label{eq:erro_generalizacao_rf}
\end{equation}

\noindent No contexto computacional finito, a predição da floresta corresponde à média aritmética simples das predições pontuais emitidas pelas $K$ árvores individuais:

\begin{equation}
    \bar{h}(\mathbf{x}) = \frac{1}{K} \sum_{k=1}^{K} h(\mathbf{x}, \Theta_k) \label{eq:predicao_rf_breiman}
\end{equation}

Embora a representação de \citeonline{Breiman_2001_RandomForests} espelhe diretamente a rotina algorítmica executada na prática, a análise assintótica de suas propriedades estatísticas demanda uma formulação analítica fechada. Com esse propósito, \citeonline{biau2012analysis} formalizou cada árvore individual como um estimador de partição linear local sobre o suporte amostral $[0, 1]^p$. Seja $A_n(\mathbf{X}; \Theta)$ a célula hiper-retangular única (folha terminal) que contém o ponto de consulta $\mathbf{X}$, a estimativa da árvore expressa-se por \cite{biau2012analysis}:

\begin{equation}
    m_n(\mathbf{X}; \Theta) = \frac{\sum_{i=1}^{n} Y_i \mathbb{I}_{\{\mathbf{X}_i \in A_n(\mathbf{X}; \Theta)\}}}{N_n(\mathbf{X}; \Theta)} \mathbb{I}_{\mathcal{E}_n(\mathbf{X}; \Theta)} \label{eq:arvore_biau}
\end{equation}

\noindent em que $N_n(\mathbf{X}; \Theta) = \sum_{i=1}^{n} \mathbb{I}_{\{\mathbf{X}_i \in A_n(\mathbf{X}; \Theta)\}}$ contabiliza o número de instâncias de treinamento contidas na partição de $\mathbf{X}$, e o evento $\mathcal{E}_n(\mathbf{X}; \Theta) = \{N_n(\mathbf{X}; \Theta) \neq 0\}$ atribui predição nula caso a célula terminal não contenha observações amostrais.

Definindo os pesos individuais $W_{n,i}(\mathbf{X}; \Theta) = \frac{\mathbb{I}_{\{\mathbf{X}_i \in A_n(\mathbf{X}; \Theta)\}}}{N_n(\mathbf{X}; \Theta)} \mathbb{I}_{\mathcal{E}_n(\mathbf{X}; \Theta)}$, a floresta agregada idealizada no limite assintótico ($K \to \infty$) resulta da esperança matemática condicional sobre a distribuição do parâmetro de aleatorização $\Theta$ \cite{biau2012analysis}:

\begin{equation}
    \bar{m}_n(\mathbf{X}) = \mathbb{E}_{\Theta}\left[m_n(\mathbf{X}; \Theta)\right] = \sum_{i=1}^{n} \mathbb{E}_{\Theta}\left[ W_{n,i}(\mathbf{X}; \Theta) \right] Y_i \label{eq:floresta_biau_pesos}
\end{equation}

A formulação (\ref{eq:floresta_biau_pesos}) evidencia que o \textit{Random Forest} opera como um estimador por ponderação local adaptativa, no qual o peso atribuído a cada observação $Y_i$ quantifica a probabilidade marginal de $\mathbf{X}_i$ compartilhar a mesma célula de particionamento com $\mathbf{X}$ ao longo do espaço induzido por $\Theta$. Ademais, \citeonline{biau2012analysis} demonstrou que esse mecanismo garante a consistência estatística do estimador e promove adaptação intrínseca à esparsidade: a taxa de convergência do erro quadrático médio depende estritamente do número de variáveis efetivamente informativas ($S \le p$), tornando o algoritmo robusto à presença de atributos ruidosos ou redundantes no espaço de entrada.

\subsubsection{Gradient Boosting}

Diferentemente do paradigma de \textit{bagging}, no qual os estimadores são construídos de maneira independente e paralela para reduzir a variância, o \textit{Gradient Boosting}, formalizado por \citeonline{friedman2001gradient}, fundamenta-se na aproximação de funções sob a perspectiva de otimização numérica no espaço funcional. Dada uma amostra de treinamento $\{\mathbf{x}_i, y_i\}_{i=1}^{N}$, o objetivo consiste em determinar uma função aproximadora $F^*(\mathbf{x})$ que minimize o valor esperado de uma função de perda diferenciável $\Psi(y, F(\mathbf{x}))$ sobre a distribuição conjunta de $(\mathbf{x}, y)$:

\begin{equation}
    F^*(\mathbf{x}) = \arg\min_{F(\mathbf{x})} \mathbb{E}_{y, \mathbf{x}}\left[\Psi(y, F(\mathbf{x}))\right] \label{eq:obj_friedman}
\end{equation}

Para tornar o problema tratável em amostras finitas, restringe-se $F(\mathbf{x})$ a uma expansão aditiva parametrizada de estimadores-base (\textit{weak learners}) $h(\mathbf{x}; \mathbf{a}_m)$, estruturados como árvores de decisão de tamanho reduzido:

\begin{equation}
    F(\mathbf{x}; \{\beta_m, \mathbf{a}_m\}_{m=0}^{M}) = \sum_{m=0}^{M} \beta_m h(\mathbf{x}; \mathbf{a}_m) \label{eq:expansao_aditiva_gb}
\end{equation}

\noindent em que $\mathbf{a}_m$ engloba os parâmetros estruturais de particionamento da $m$-ésima árvore (variáveis de divisão, pontos de corte e médias dos nós terminais) e $\beta_m$ denota o coeficiente de expansão.

Partindo de uma estimativa inicial constante $F_0(\mathbf{x}) = \arg\min_{\rho} \sum_{i=1}^{N} \Psi(y_i, \rho)$, o algoritmo adota uma estratégia gulosa passo a passo análoga ao método de descida mais íngreme. Em cada iteração $m \in \{1, \dots, M\}$, calcula-se a direção do gradiente negativo não restrito da função de perda em relação ao modelo acumulado $F_{m-1}(\mathbf{x})$, originando as chamadas pseudo-respostas (ou pseudo-resíduos) $\tilde{y}_i$:

\begin{equation}
    \tilde{y}_i = -\left[\frac{\partial \Psi(y_i, F(\mathbf{x}_i))}{\partial F(\mathbf{x}_i)}\right]_{F(\mathbf{x}) = F_{m-1}(\mathbf{x})}, \quad \forall i \in \{1, \dots, N\} \label{eq:pseudo_residuos}
\end{equation}

No caso particular de regressão por mínimos quadrados ordinários, em que $\Psi(y, F) = \frac{1}{2}(y - F)^2$, a Equação (\ref{eq:pseudo_residuos}) reduz-se diretamente aos resíduos correntes do comitê, $\tilde{y}_i = y_i - F_{m-1}(\mathbf{x}_i)$. Em seguida, ajusta-se uma nova árvore de decisão $h(\mathbf{x}; \mathbf{a}_m)$ por mínimos quadrados para aproximar a direção dos pseudo-resíduos e determina-se o comprimento ótimo do passo $\rho_m$ via busca linear:

\begin{align}
    \mathbf{a}_m &= \arg\min_{\mathbf{a}, \beta} \sum_{i=1}^{N} \left[\tilde{y}_i - \beta h(\mathbf{x}_i; \mathbf{a})\right]^2 \label{eq:ajuste_arvore_gb} \\
    \rho_m &= \arg\min_{\rho} \sum_{i=1}^{N} \Psi\left(y_i, F_{m-1}(\mathbf{x}_i) + \rho h(\mathbf{x}_i; \mathbf{a}_m)\right) \label{eq:line_search_gb}
\end{align}

Para mitigar o \textit{overfitting} decorrente da adição sucessiva de árvores, \citeonline{friedman2001gradient} introduz um mecanismo de regularização por encolhimento proporcional (\textit{shrinkage}), escalonando a contribuição de cada novo estimador por um parâmetro $0 < \nu \le 1$ (taxa de aprendizado), o que resulta na regra de atualização iterativa:

\begin{equation}
    F_m(\mathbf{x}) = F_{m-1}(\mathbf{x}) + \nu \rho_m h(\mathbf{x}; \mathbf{a}_m), \quad 0 < \nu \le 1 \label{eq:update_shrinkage_gb}
\end{equation}

\noindent estabelecendo um compromisso direto entre a magnitude de $\nu$ e o número total de iterações $M$ necessárias para a convergência do comitê.

A partir do arcabouço clássico de \citeonline{friedman2001gradient}, implementações modernas como o XGBoost (\textit{eXtreme Gradient Boosting}) e o LightGBM (\textit{Light Gradient Boosting Machine}) introduziram aprimoramentos analíticos e computacionais voltados à regularização estrutural e à escalabilidade.

Proposto por \citeonline{Chen_2016_XGBoost}, o XGBoost define o espaço de árvores de regressão \gls{CART} como $\mathcal{F} = \{f(\mathbf{x}) = w_{q(\mathbf{x})}\}$, em que a estrutura $q: \mathbb{R}^p \to \{1, \dots, T\}$ mapeia cada observação ao índice de seu respectivo nó terminal e $\mathbf{w} \in \mathbb{R}^T$ representa o vetor de pesos associados às $T$ folhas da árvore. O algoritmo distingue-se da formulação tradicional de primeira ordem pela incorporação de um termo de penalização explícito sobre a complexidade do modelo e pelo emprego de uma aproximação em série de Taylor de segunda ordem. Na $t$-ésima iteração, busca-se o estimador $f_t \in \mathcal{F}$ que minimiza o objetivo regularizado:

\begin{equation}
    \mathcal{L}^{(t)} = \sum_{i=1}^{N} l\left(y_i, \hat{y}_i^{(t-1)} + f_t(\mathbf{x}_i)\right) + \Omega(f_t), \quad \text{com} \quad \Omega(f_t) = \gamma T + \frac{1}{2}\lambda \|\mathbf{w}\|^2 \label{eq:obj_xgboost}
\end{equation}

\noindent em que $l$ é uma função de perda convexa diferenciável que calcula a diferença entre o valor alvo $y_i$ e a predição $\hat{y}_i$, $\gamma$ penaliza a quantidade de folhas $T$ e $\lambda$ controla a regularização $L_2$ sobre a magnitude dos pesos terminais. Expandindo a função de perda em série de Taylor até a segunda ordem em torno de $\hat{y}^{(t-1)}$ e removendo os termos constantes, obtém-se o objetivo simplificado:

\begin{equation}
    \tilde{\mathcal{L}}^{(t)} = \sum_{i=1}^{N} \left[ g_i f_t(\mathbf{x}_i) + \frac{1}{2} h_i f_t^2(\mathbf{x}_i) \right] + \gamma T + \frac{1}{2}\lambda \|\mathbf{w}\|^2 \label{eq:taylor_xgboost}
\end{equation}

\noindent no qual $g_i = \partial_{\hat{y}^{(t-1)}} l(y_i, \hat{y}^{(t-1)})$ e $h_i = \partial^2_{\hat{y}^{(t-1)}} l(y_i, \hat{y}^{(t-1)})$ denotam as estatísticas de gradiente de primeira e segunda ordem, respectivamente. Definindo $I_j = \{i \mid q(\mathbf{x}_i) = j\}$ como o conjunto de instâncias alocadas na folha $j$, calcula-se analiticamente o peso ótimo $w_j^*$ e o escore de qualidade estrutural correspondente:

\begin{equation}
    w_j^* = -\frac{\sum_{i \in I_j} g_i}{\sum_{i \in I_j} h_i + \lambda}, \qquad \tilde{\mathcal{L}}^{(m)}(q) = -\frac{1}{2} \sum_{j=1}^{T} \frac{\left(\sum_{i \in I_j} g_i\right)^2}{\sum_{i \in I_j} h_i + \lambda} + \gamma T \label{eq:pesos_otimos_xgboost}
\end{equation}

\noindent A partir da Equação (\ref{eq:pesos_otimos_xgboost}), a seleção gulosa de pontos de corte avalia a redução da perda $\mathcal{L}_{\text{split}}$ resultante da partição de um nó $I = I_L \cup I_R$ em subnós esquerdo ($I_L$) e direito ($I_R$):

\begin{equation}
    \mathcal{L}_{\text{split}} = \frac{1}{2} \left[ \frac{\left(\sum_{i \in I_L} g_i\right)^2}{\sum_{i \in I_L} h_i + \lambda} + \frac{\left(\sum_{i \in I_R} g_i\right)^2}{\sum_{i \in I_R} h_i + \lambda} - \frac{\left(\sum_{i \in I} g_i\right)^2}{\sum_{i \in I} h_i + \lambda} \right] - \gamma \label{eq:split_xgboost}
\end{equation}

\noindent Adicionalmente ao objetivo regularizado e ao encolhimento por taxa de aprendizado, o XGBoost incorpora a subamostragem aleatória de atributos inspirada no Random Forest, algoritmos conscientes de esparsidade e esboços de quantis ponderados pelas estatísticas de segunda ordem $h_i$ para discretização eficiente de atributos contínuos.

Embora o \textit{XGBoost} mitigue o custo de ordenação por meio de estruturas de blocos em memória e quantis ponderados, a complexidade computacional para construir os histogramas e avaliar a Equação (\ref{eq:split_xgboost}) permanece proporcional tanto ao número total de observações $N$ quanto à dimensionalidade de atributos $p$. Visando superar esse gargalo em bases de dados de larga escala sem comprometer a acurácia preditiva, \citeonline{ke2017lightgbm} propuseram o algoritmo LightGBM.

O LightGBM preserva integralmente a expansão de Taylor de segunda ordem e a função objetivo regularizada expressas nas Equações (\ref{eq:taylor_xgboost} - \ref{eq:split_xgboost}), mas introduz três inovações estruturais e algorítmicas:

\begin{enumerate}
    \item \textbf{Particionamento Baseado em Histogramas e Crescimento \textit{Leaf-wise}:} Os atributos contínuos são discretizados previamente em $k$ compartimentos inteiros, reduzindo a complexidade de busca do ponto de corte de $\mathcal{O}(N \cdot p)$ para $\mathcal{O}(k \cdot p)$, além de permitir a obtenção do histograma de um nó irmão por simples subtração vetorial em relação ao nó pai. Adicionalmente, em contraste com a estratégia tradicional de expansão por nível de profundidade, o LightGBM adota o crescimento orientado às folhas, particionando iterativamente apenas a folha terminal que proporciona o maior ganho global $\mathcal{L}_{\text{split}}$ em toda a árvore, sujeito a um limite máximo de profundidade para evitar sobreajuste;

    \item \textbf{\gls{GOSS}:} Fundamenta-se na premissa teórica de que instâncias associadas a gradientes de primeira ordem $|g_i|$ de pequena magnitude já se encontram bem ajustadas pelo comitê corrente, contribuindo marginalmente para o ganho de informação. Em cada iteração, o método \gls{GOSS} ordena as observações de forma decrescente por $|g_i|$, retém integralmente o subconjunto $A$ correspondente à fração $a \times 100\%$ de instâncias com os maiores gradientes e sorteia aleatoriamente um subconjunto $B$ de tamanho $b \times |A^c|$ dentre o complemento de instâncias com gradientes reduzidos. Para preservar a distribuição original da perda ao avaliar os candidatos a ponto de corte $d$ sobre o atributo $j$, os gradientes pertencentes a $B$ são amplificados pelo fator constante $\frac{1-a}{b}$, resultando no estimador de ganho aproximado $\tilde{V}_j(d)$:
    
    \begin{equation}
        \tilde{V}_j(d) = \frac{1}{N} \left[ \frac{\left(\sum_{x_i \in A_l} g_i + \frac{1-a}{b}\sum_{x_i \in B_l} g_i\right)^2}{n_l^j(d)} + \frac{\left(\sum_{x_i \in A_r} g_i + \frac{1-a}{b}\sum_{x_i \in B_r} g_i\right)^2}{n_r^j(d)} \right] \label{eq:goss_lightgbm}
    \end{equation}
    
    \noindent em que $A_l = \{x_i \in A \mid x_{ij} \le d\}$, $A_r = \{x_i \in A \mid x_{ij} > d\}$, $B_l = \{x_i \in B \mid x_{ij} \le d\}$, $B_r = \{x_i \in B \mid x_{ij} > d\}$, e os denominadores $n_l^j(d)$ e $n_r^j(d)$ contabilizam o total de observações alocadas aos subnós esquerdo e direito, respectivamente;

    \item \textbf{\gls{EFB}:} Em espaços de alta dimensão com presença de esparsidade, múltiplos atributos são mutuamente exclusivos, isto é, raramente assumem valores não nulos de maneira simultânea. O método EFB reduz a quantidade efetiva de atributos ao fundir variáveis mutuamente exclusivas em um único feixe, deslocando os intervalos de seus \textit{bins} originais por constantes de compensação para preservar a identificabilidade de cada variável. A determinação dos feixes ótimos é reduzida a um problema de coloração de grafos \textit{NP-difícil} — no qual os vértices representam os atributos e as arestas ponderam os conflitos de não exclusividade —, resolvido eficientemente por uma heurística gulosa com taxa de conflito máxima tolerada.
\end{enumerate}

\section{Métricas}

Por fim, há também a necessidade de se otimizar as métricas de cada um dos modelos, visto que esses determinarão o quanto o modelo se adequa ao problema, e o quanto ele consegue predizer os valores da variável objetivo. Sendo assim, abaixo serão apresentadas as métricas utilizadas e que foram otimizadas:

\begin{itemize}
    \item \textbf{Coeficiente de Determinação ($\text{R}^2$):} O $\text{R}^2$ mede a proporção da variância total na variável alvo que é explicada pelo modelo. É a métrica primária de precisão.
    \begin{equation}
        R^2 = 1 - \frac{\sum_{i=1}^{n} (y_i - \hat{y}_i)^2}{\sum_{i=1}^{n} (y_i - \bar{y})^2} = 1 - \frac{\text{RSS}}{\text{TSS}}
    \end{equation}
    Onde \gls{RSS} é a Soma dos Quadrados dos Resíduos (erro do modelo) e \gls{TSS} é a variância total dos dados. Um $\text{R}^2$ próximo de 1 indica um ajuste excelente.
    
    \item \textbf{Raiz do Erro Quadrático Médio (RMSE):} O \gls{RMSE} é a principal métrica de erro, representando a magnitude média dos desvios. É a raiz quadrada da média dos erros quadrados, mantendo a unidade de medida original (bbl/dia) para interpretação prática.
    \begin{equation}
        \text{\gls{RMSE}} = \sqrt{\frac{1}{n} \sum_{i=1}^{n} (y_i - \hat{y}_i)^2}
    \end{equation}
    A função quadrática dentro da raiz penaliza erros grandes de forma desproporcional, tornando o \gls{RMSE} sensível a \textit{outliers}.

    \item \textbf{Erro Absoluto Médio (MAE):} O \gls{MAE} mede a média dos erros em valor absoluto.
    \begin{equation}
         \text{\gls{MAE}} = \frac{1}{n} \sum_{i=1}^{n} |y_i - \hat{y}_i|
    \end{equation}
    
    Por não elevar o erro ao quadrado, o MAE é mais robusto contra \textit{outliers} do que o \gls{RMSE} e fornece uma estimativa de erro mais intuitiva.

    \item \textbf{Diferença do Coeficiente de Determinação ($\Delta \text{R}^2$):} Esta métrica avalia a capacidade de \textbf{generalização} do modelo, sendo essencial para identificar o \textit{overfitting}.
    \begin{equation}
        \Delta \text{R}^2 = \text{R}^2_{\text{Treino}} - \text{R}^2_{\text{Teste}}
    \end{equation}
    Um $\Delta \text{R}^2$ próximo de zero indica que o desempenho do modelo no conjunto de teste é similar ao do treino (boa generalização). Um valor positivo alto é um forte indicador de que o modelo memorizou os dados de treino, perdendo poder preditivo em dados não vistos.
\end{itemize}